\chapter{Taxonomy over QUIC Datapaths}
\label{chap:taxonomy}

To comprehensively analyze QUIC datapaths, we compare \ac{DPDK}, in-kernel QUIC (lxin-quic), and AF\_XDP, describing the unique tradeoffs inherent in each approach.

\section{Implementation Complexity}
\label{subsec:implementation-complexity}

An immediate benefit of lxin-quic is its proposed seamless integration into the standardized Linux socket APIs. This means applications already using the Linux socket API for \ac{TCP} can leverage the same networking system calls such as \texttt{listen(2)}, \texttt{accept(2)}, \texttt{connect(2)}, \texttt{sendmsg(2)}, and \texttt{recvmsg(2)} for QUIC\footnote{https://lwn.net/Articles/1028932/}. This significantly reduces the friction associated with adopting QUIC. In contrast, userspace QUIC implementations lack a standardized API, requiring a distinct integration process for each. Consequently, \ac{DPDK} and AF\_XDP, which primarily support userspace QUIC, become inherently more complex to use.

Moreover, there is also the concept of availability. \ac{DPDK}, for instance, is limited to a specific set of \ac{NIC}s and drivers\footnote{https://core.dpdk.org/supported/\#nics}. While XDP\_SKB mode enjoys native Linux support, its XDP\_DRV counterpart requires specific driver implementation, though support for this has been increasing with each new kernel version\footnote{https://docs.ebpf.io/linux/program-type/BPF\_PROG\_TYPE\_XDP/\#driver-support}. Lxin-quic, as a loadable kernel module, must be built against a precise kernel version due to the unstable kernel ABI across releases. While in-kernel QUIC might eventually be integrated into the core kernel codebase, its current modular status places an additional maintenance burden on its developers.

\section{Performance Improvement Potential}
\label{subsec:performance-improvement-potential}

One niche application for transport protocols is serving files from storage over the network, similar what a streaming service provider like Netflix would do. As shown in Figure~\ref{fig:memory-bandwidth-QUIC}, the problem with userspace QUIC is that it may not use the \texttt{sendfile(2)} system call, which is beneficial due to avoiding a copy into userspace. For this use case, in-kernel QUIC is advantageous. For \ac{DPDK}, while not having access to \texttt{sendfile(2)}, it may use other methods such as \ac{SPDK}\footnote{https://spdk.io/} to achieve high goodput. Finally, AF\_XDP does not offer any advantage as it is simply a datapath from the kernel to userspace within the context of networking.

\section{Unwanted Side Effects}
\label{subsec:unwanted-side-effects}

One significant drawback of \ac{DPDK} stems from its kernel bypass approach. Here, \ac{DPDK} forfeits the benefits of the established Linux kernel networking stack. Consequently, this requires recreating a complete network stack, including drivers, within the \ac{DPDK} domain. Furthermore, the entire ecosystem of tools for monitoring, managing, and troubleshooting \ac{DPDK} applications also needs to be duplicated or developed from scratch. \ac{DPDK}'s smaller user base compared to the Linux kernel also leads to higher maintenance costs as the maintencance burden cannot be spread across more developers. Finally, \ac{DPDK} gains exclusive control over a \ac{NIC} by detaching it from the kernel, rendering the \ac{NIC} 'invisible' to standard Linux networking tools. This becomes particularly problematic with single-\ac{NIC} systems, as it prevents other userspace applications from accessing the network.

While moving QUIC into the Linux kernel could reduce packet I/O overhead, a fundamental challenge, as discussed in Subsection~\ref{subsec:slow-adoption-of-new-features}, is the inherent slowdown in development velocity for kernel software. For instance, a bug fix necessitates either patching and recompiling the kernel or upgrading the entire kernel version, affecting every application on the operating system. In contrast, userspace applications can independently manage their versions. Debugging kernel-level QUIC also remains a complex endeavor. Options like print debugging, \ac{eBPF} programs, or perf are available, but comprehensive step-through debugging often requires emulating the entire system with tools like QEMU\footnote{https://www.qemu.org/}. This adds significantly more friction compared to debugging in userspace.

A notable disadvantage of AF\_XDP is its requirement for in-depth knowledge of kernel internals, including \ac{eBPF}, \ac{XDP} hooks, and AF\_XDP-specific features like ring buffers and the UMEM region. Moreover, despite ongoing improvements, documentation remains sparse, especially concerning driver-related issues, as detailed in Subsection~\ref{subsec:xdp-on-the-client}. This contrasts with in-kernel QUIC, which would leverage standardized networking APIs already familiar to many developers.

\section{Bring-up and Privilege Level Usage}
\label{subsec:bring-up-and-privelege-level-usage}

Bringing up lxin-quic necessitates loading it as a kernel module, which, like other kernel modules, requires root privileges or the \url{CAP\_SYS\_MODULE} capability\footnote{https://man7.org/linux/man-pages/man7/capabilities.7.html}. Similarly, \ac{DPDK} depends on the \texttt{vfio-pci} kernel module for \ac{NIC} interaction and thus also requires higher privileges.

Regarding \ac{eBPF}, while certain limited program types like \url{BPF\_PROG\_TYPE\_SOCKET\_FILTER} can be loaded without elevated privileges, \ac{XDP} programs (of type \url{BPF\_PROG\_TYPE\_XDP}) require higher privileges\footnote{https://man7.org/linux/man-pages/man2/bpf.2.html}. However, a key advantage of \ac{XDP}, which is particularly beneficial for applications such as \ac{DDoS} firewalls, is its ability to atomically swap \ac{eBPF} programs at runtime\footnote{https://lwn.net/Articles/815707/}. This avoids network traffic interference and allows the \ac{XDP} program's logic to adapt dynamically to changing traffic patterns.